% Defesa de TCC em Computação com o tema metropolis
% Modelo do ÉguaTex (https://eguatex.github.io), em domínio público pela licença CC0 1.0:
% use, modifique e distribua sem pedir permissão e sem citar a origem.
%
% Compile com pdfLaTeX e BibTeX (o ÉguaTex faz isso sozinho).
% Arquivos: main.tex (os slides) e referencias.bib (bibliografia).
%
% O metropolis foi feito para XeLaTeX, mas aqui ele roda no pdfLaTeX: as partes de
% layout e cor do tema entram uma a uma, e a Fira Sans vem do pacote FiraSans.

\documentclass[11pt,aspectratio=169]{beamer}

% ---------- idioma e fontes ----------
\usepackage[T1]{fontenc}
\usepackage[utf8]{inputenc}
\usepackage[brazilian]{babel}
\usepackage[sfdefault,book]{FiraSans}
\usepackage[scale=0.88]{FiraMono}

% ---------- tema metropolis sem a parte de fontes (que exige XeLaTeX) ----------
\useinnertheme{metropolis}
\useoutertheme[progressbar=frametitle, numbering=fraction]{metropolis}
\usecolortheme{metropolis}
\setbeamerfont{title}{size=\Large, series=\bfseries}
\setbeamerfont{frametitle}{size=\large}
\setbeamerfont{section title}{size=\Large}
\setbeamerfont{standout}{size=\Large, series=\bfseries}
\setbeamerfont{block title}{series=\bfseries}
\setbeamerfont{page number in head/foot}{size=\scriptsize}
\setbeamertemplate{navigation symbols}{}
\setbeamertemplate{bibliography item}{\insertbiblabel}

% ---------- código, diagramas e gráficos ----------
\usepackage{listings}
\usepackage{tikz}
\usetikzlibrary{arrows.meta, positioning}
\usepackage{pgfplots}
\pgfplotsset{compat=1.18}
\usepackage{booktabs}

% Estilo dos trechos de código. Dentro de lstlisting, o texto entre | e |
% é LaTeX: \alert<2>{...} destaca uma parte só no passo 2 do slide.
\lstset{
  basicstyle=\ttfamily\small,
  keywordstyle=\color{mLightBrown}\bfseries,
  commentstyle=\color{mDarkTeal!60},
  stringstyle=\color{mLightGreen},
  numbers=left, numberstyle=\tiny\color{gray}, numbersep=6pt,
  columns=fullflexible, keepspaces=true,
  escapechar=|,
  language=Python,
  showstringspaces=false,
}

% ---------- instruções de preenchimento ----------
% As notas em cinza explicam o que pôr em cada slide.
% Troque \instrucoestrue por \instrucoesfalse para esconder todas de uma vez.
\newif\ifinstrucoes
\instrucoestrue
\newcommand{\instrucao}[1]{\ifinstrucoes\par\vfill{\color{gray}\scriptsize\itshape #1\par}\fi}

% ---------- dados da capa: preencha aqui ----------
\title{Detecção de anomalias em registros de servidores com aprendizado de máquina}
\subtitle{Defesa do Trabalho de Conclusão de Curso}
\author{Nome do autor \quad{\small(orientação: Nome do orientador)}}
\institute{Bacharelado em Ciência da Computação, Universidade Federal Exemplo}
\date{Cidade, Ano}

\begin{document}

\maketitle

% ======================================================================
\section{Motivação}
% ======================================================================
\begin{frame}{O problema}
  \begin{itemize}
    \item Um servidor web gera milhões de linhas de registro por dia.
    \item Falhas raras ficam escondidas no meio das mensagens normais.
    \item \alert{Pergunta:} um modelo não supervisionado encontra essas falhas a tempo?
  \end{itemize}
  \instrucao{Mostre o problema com um número concreto e feche com a pergunta de pesquisa. Três tópicos bastam.}
\end{frame}

% ======================================================================
\section{Proposta}
% ======================================================================
\begin{frame}{Arquitetura da solução}
  \centering
  \begin{tikzpicture}[
      etapa/.style={draw=mDarkTeal, fill=mDarkTeal!8, rounded corners=3pt,
        minimum width=2.4cm, minimum height=1cm, align=center, font=\small},
      seta/.style={-{Stealth[length=6pt]}, thick, mDarkTeal},
      node distance=0.7cm]
    \node[etapa] (coleta) {Coleta\\dos registros};
    \node[etapa, right=of coleta] (parse) {Extração\\de padrões};
    \node[etapa, right=of parse] (vetor) {Vetores de\\contagem};
    \node[etapa, right=of vetor, fill=mLightBrown!20, draw=mLightBrown] (modelo) {Modelo de\\anomalias};
    \draw[seta] (coleta) -- (parse);
    \draw[seta] (parse) -- (vetor);
    \draw[seta] (vetor) -- (modelo);
    \node[below=0.5cm of modelo, font=\small, text=mLightBrown] {alerta ao operador};
    \draw[seta, mLightBrown] (modelo) -- ++(0,-0.75);
  \end{tikzpicture}
  \instrucao{Desenhe o fluxo da solução em no máximo cinco caixas. Destaque com outra cor a etapa que é contribuição sua.}
\end{frame}

% ======================================================================
\section{Implementação}
% ======================================================================
\begin{frame}[fragile]{O laço principal, passo a passo}
\begin{lstlisting}
def detectar(janelas, modelo, limiar):
    alertas = []
    for janela in janelas:
        |\alert<2>{vetor = contar\_eventos(janela)}|
        |\alert<3>{nota = modelo.pontuar(vetor)}|
        if |\alert<4>{nota > limiar}|:
            alertas.append(janela.inicio)
    return alertas
\end{lstlisting}
  \begin{itemize}
    \item<2-> Cada janela de 60 s vira um vetor de contagens.
    \item<3-> O modelo devolve uma nota de anomalia.
    \item<4-> O limiar controla o equilíbrio entre falsos alarmes e falhas perdidas.
  \end{itemize}
  \instrucao{Mostre só o trecho que explica a ideia, com até 10 linhas. Cada passo do slide destaca uma linha e acrescenta a explicação.}
\end{frame}

% ======================================================================
\section{Avaliação}
% ======================================================================
\begin{frame}{Resultados}
  \begin{columns}[c]
    \begin{column}{0.55\textwidth}
      \begin{tikzpicture}
        \begin{axis}[
          width=\linewidth, height=5cm,
          xlabel={limiar}, ylabel={medida},
          xmin=0.5, xmax=0.95, ymin=0, ymax=1,
          legend pos=south west, legend style={font=\scriptsize},
          tick label style={font=\scriptsize}, label style={font=\small},
          /pgf/number format/use comma,
        ]
          \addplot+[thick, mark=*] coordinates {(0.5,0.42) (0.6,0.55) (0.7,0.71) (0.8,0.83) (0.9,0.91) (0.95,0.94)};
          \addlegendentry{precisão}
          \addplot+[thick, mark=square*] coordinates {(0.5,0.97) (0.6,0.93) (0.7,0.88) (0.8,0.79) (0.9,0.62) (0.95,0.48)};
          \addlegendentry{revocação}
        \end{axis}
      \end{tikzpicture}
    \end{column}
    \begin{column}{0.4\textwidth}
      \begin{table}
        \small
        \begin{tabular}{@{}lr@{}}
          \toprule
          Métrica & Valor \\
          \midrule
          Precisão & 0,83 \\
          Revocação & 0,79 \\
          F1 & 0,81 \\
          \bottomrule
        \end{tabular}
        \par\smallskip{\scriptsize Limiar 0,8 (dados de exemplo).\par}
      \end{table}
    \end{column}
  \end{columns}
  \instrucao{Compare com uma linha de base e diga a configuração usada. Troque os pontos de exemplo pelos seus resultados.}
\end{frame}

% ======================================================================
\section{Conclusão}
% ======================================================================
\begin{frame}{Conclusões}
  \begin{itemize}
    \item O modelo encontrou 79\% das falhas com 17\% de falsos alarmes.
    \item A extração de padrões reduziu os dados a 1\% do volume original.
    \item \alert{Trabalho futuro:} treinar o modelo em fluxo contínuo.
  \end{itemize}
  \instrucao{Responda à pergunta de pesquisa com números, diga uma limitação e um próximo passo.}
\end{frame}

\begin{frame}[standout]
  Obrigado!\\[6pt]
  {\normalsize Perguntas?}
\end{frame}

\begin{frame}{Referências}
  \nocite{cormen2009, knuth1998}
  \small
  \bibliographystyle{plain}
  \bibliography{referencias}
\end{frame}

\end{document}
